%%===================================================================%%
%% Discussion and Conclusions — first draft, 2026-09-21               %%
%% Goes between \section{Results} and \section{Methods}.              %%
%% Every number here also appears in results_body_v2.tex; nothing     %%
%% new is computed in this section.                                   %%
%%===================================================================%%

\section{Discussion}

\subsection{What the temporal split costs, and what the cost is made of}

Both benchmarks lose discrimination under a 2010 temporal split: AUROC falls by
0.088 in DILIrank and 0.077 in DICTrank against a matched random reference, with
permutation $p$ of 0.0070 and 0.0050. Numbers of that size are the usual finding
of a time-split study, and on their own they would not warrant a paper. What the
AUROC does not show is that the same split raises expected calibration error from
0.121 to 0.438 and from 0.126 to 0.270, and that the Murphy decomposition assigns
86.1\% and 82.5\% of the increase in Brier score to the reliability term rather
than the resolution term. (The random-split figures quoted here come from a single
matched partition; the 200-partition controls of the next section give 0.131 and
0.130 for the same quantity, and the two differ only in the second decimal place.)

The practical difference between the two failures is that discrimination is a
property of the ranking and calibration is a property of the number. A model that
loses 0.09 of AUROC still ranks most compounds sensibly. A model whose mean
predicted probability is 0.732 on a held-out set with a 0.308 positive rate is
telling a user something false about every compound it scores. In this setting the
number is what gets acted on: a predicted probability of hepatotoxicity is read as
a probability, compared against a threshold, and used to decide whether a compound
advances. Seventy-one of 156 DILIrank held-out compounds carried a predicted
probability above 0.9, at a mean of 0.962, and 40.8\% of them were positive.

The failure is systematic rather than random, which is why the optimal decision
threshold moves to 0.920 and 0.940 rather than staying near 0.5, and why moving
it recovers a substantial part of the classification performance lost at the
default. A threshold that high is not a usable default; the point is the distance
one has to travel to find it, and the fact that nothing in the AUROC signals that
the journey is necessary.

\subsection{Scaffold splitting is not a proxy for temporal validation}

Structural novelty is the usual stand-in for prospective difficulty, on the
reasoning that future compounds are structurally new. In these two datasets the
substitution does not hold.

Measured by AUROC, Bemis--Murcko scaffold cross-validation is not reliably harder
than stratified cross-validation: it raised DILIrank by 0.002 and lowered DICTrank
by 0.022, both within a fold-to-fold standard deviation, and in opposite
directions. Measured by ECE, scaffold and stratified cross-validation agree in
both datasets to within a standard deviation, while the temporal split on the same
compounds gives ECE 3.4 and 2.0 times larger.

The mechanism is visible in the scaffold counts. Of 542 unique scaffolds among the
796 DILIrank compounds, 455 (84\%) appear exactly once; of 638 among the 919
DICTrank compounds, 545 (85\%). Grouping by scaffold cannot separate what is
already separated: when most scaffolds are singletons, a scaffold-grouped fold is
close to a random fold, and the design converges on the thing it was meant to be
harder than. This is a property of drug-like reference sets rather than of these
two datasets, and it should be checked and reported whenever scaffold splitting is
offered as evidence of prospective validity. The check is one line of code and the
number belongs in the Methods of any paper that uses the design.

None of this argues that scaffold splitting is uninformative. It argues that
structural novelty and temporal novelty are different axes, that these datasets
separate them, and that a model can be robust on one while failing on the other.

\subsection{What does explain the collapse, and what does not}

Three candidate explanations were tested directly. One caveat applies throughout:
only the prevalence controls were evaluated against calibration. The chemical-space
and label-maturity analyses measure the discrimination gap alone, so they can rule
an explanation in or out for AUROC but not for ECE.

\textbf{Prevalence shift is a large part of it.} The positive rate falls from
65.5\% to 30.8\% across the DILIrank boundary and from 77.0\% to 61.4\% across the
DICTrank one, and a model calibrated to one base rate is miscalibrated on another
by construction. Matching both base rates in a random split closes 62\% and 66\%
of the ECE excess. Both datasets agree closely on that fraction, which is itself
informative: the effect is not idiosyncratic. But 38\% and 34\% of the excess survive
full prevalence matching, and the residual remains significant after
correction ($q \le 0.0348$). Prevalence is a major term, not the whole term.

Notably, prevalence does not explain any of the \emph{discrimination} loss. Matching
the held-out positive rate to the training rate moved DILIrank AUROC from 0.648 to
0.632, widening the gap rather than closing it. The two failures have different
causes, which is a further reason not to treat AUROC and ECE as interchangeable
summaries of the same phenomenon.

\textbf{Chemical-space novelty is not supported as an explanation of the
discrimination gap.} Held-out
compounds do sit far from the training set --- median nearest-neighbour Tanimoto
0.271, with 73.7\% below 0.4 --- and AUROC rises across the three similarity bins
large enough to evaluate, from 0.543 to 0.711 to 0.857; the highest-similarity bin
holds 12 compounds and was not evaluated. But splitting the held-out set at the
median gives a difference of only
0.068 at a bootstrap $p$ of 0.527, and the correlation between similarity and the
predicted probability of the true class is weak ($\rho = 0.116$, $p = 0.148$).
Neither survives correction. The bin-wise trend is consistent with a real effect
too small to detect at $n = 156$, and it is reported as such rather than as a
finding. As with label maturity, calibration was not recomputed by similarity
stratum, so this analysis says nothing about whether structural distance
contributes to the calibration collapse.

\textbf{Label maturity contributes to the discrimination gap in one dataset and not
the other.} Post-marketing labels accumulate, so young compounds are
under-labelled, and the gradient is present: DILIrank positive rate rises from
14.0\% at seven years or less to 67.7\% at 16--25 years. Restricting the DILIrank
held-out set to compounds with at least twelve years of observation narrows the
AUROC gap by 24\%, but only on 50 compounds, and down-weighting negatives by
observation time changes it by less than 0.003. DICTrank does not move at all: the
most favourable control reduces the gap by 2\%, and the eight-year restriction
widens it. In both datasets the comparison that holds training and held-out
maturity to the same standard --- the only one that is a fair test --- widens the
gap rather than narrowing it.

Two further observations bear on how much weight the maturity story can carry. The
gradient does not continue: in the oldest stratum, which holds the majority of both
datasets, the positive rate falls back, so label accumulation appears to saturate
rather than to continue indefinitely. And the maturity controls were evaluated on
AUROC alone; calibration was not measured under them. Whether label censoring also
contributes to the calibration collapse is therefore untested here, and no claim is
made about it either way. The evidence that prevalence shift is the dominant
calibration term comes from the nested controls of the previous subsection, which
are independent of this one.

\subsection{Recalibration is not a remedy, and the reason matters}

An oracle Platt scaling fitted on the held-out labels reduces ECE to 0.048 and
0.023, so the miscalibration is in principle correctable by a monotone map. Under
a prospective constraint it is not corrected. Internal Platt scaling and isotonic
regression, fitted on the most recent quarter of the training compounds, close
8--10\% of the available gap in DILIrank and 40--42\% in DICTrank.

That the two datasets differ by a factor of four to five is the interesting part.
An internal calibration set fixes the miscalibration a model would have had at the
end of the training period. It cannot anticipate a shift that happens afterwards.
One reading of the difference is that more of the shift falls beyond the last
training year in DILIrank, whose base rate moves by 35 percentage points across the
boundary, than in DICTrank, where it moves by 16. That is a hypothesis consistent
with two datasets, not a tested mechanism, and it would take more endpoints to
distinguish it from other explanations.

A correction that explicitly rescales for a change in base rate recovers 38\% in
DILIrank and 12\% in DICTrank, the reverse ordering, and for the same reason. But
that correction is not prospective: it uses the held-out positive rate, which a
user would have to estimate rather than know. It is reported because it isolates
the base-rate component, not as a method anyone can deploy. Estimating a future
prevalence of drug-induced liver injury among not-yet-approved compounds is not
obviously easier than the prediction problem itself.

The practical conclusion is narrow and, we think, defensible: recalibration on
historical data is worth doing and is not sufficient, and a model card that
reports only internally-calibrated performance is reporting the easy half of the
problem.

\subsection{An assumption-free confirmation}

Split conformal prediction provides a coverage guarantee that depends on
exchangeability and on nothing else --- no distributional assumption, no
calibration metric, no binning choice. A temporal split violates exchangeability
by construction, so the guarantee is expected to fail, and it does: marginal
coverage at a nominal 90\% falls to 73.7\% and 82.4\%, in the same order as the
ECE ratios.

The class-conditional decomposition adds something the marginal figure hides. The
loss falls almost entirely on the negative class --- 63.9\% and 55.3\% against
nominal 90\%, while the positive class stays at or above nominal. In this domain
the negative label is the reassuring one, so a prediction set that omits it is a
model declining to certify as safe a compound the label calls safe. Marginal
coverage reports a guarantee mildly violated; the decomposition shows that one
half of it has gone and the other has not. Split conformal guarantees only
marginal coverage, so these class-conditional numbers are descriptive rather than
tested, which is why the random reference itself can fall below nominal within a
class.

\subsection{Relation to previous work}

Time-split validation has been standard since Sheridan, and has been applied to
DILI prediction before. Calibration degradation under distribution shift has also
been reported: Morger and colleagues documented it across ChEMBL endpoints and
showed that replacing the calibration set mitigates it, and Friesacher and
colleagues showed that temporal shift in industrial pharmaceutical data degrades
uncertainty quantification. Crucitti and Mosquera Orgueira recently argued that the
choice of train--test split matters more for evaluation than for the measured
performance itself, which is the same observation this work makes from a different
direction.

What this study adds is the combination of three things. The endpoints are
regulatory labels for clinical toxicity rather than biochemical activity, and the
time axis is the drug approval year rather than a publication date or a database
release; these are the quantities a regulatory or portfolio decision actually turns
on. The prevalence confound is separated from the rest with nested controls rather
than assumed away. And the calibration claim is made three times over --- through
ECE and the Murphy decomposition, through a distribution-free conformal argument,
and through sensitivity analyses over seeds, binning schemes, cut-off years,
learners and representations --- so that it does not rest on any one metric's
known weaknesses.

\subsection{Limitations}

The held-out sets are small. DILIrank contributes 156 compounds with 48 positives
and a bootstrap AUROC interval of 0.555--0.736. Subgroup cells range from 10 to 204
compounds and none of the twelve tests survives multiple-comparison correction; the
below-chance AUROC in the 32-compound antineoplastic subgroup is reported for
completeness and supports no claim here.

The approval year is a proxy for the date at which a compound entered the world,
not for the date at which its toxicity label was fixed. An audit identified seven
compounds whose earliest matching approval record predates the split boundary, and
three DICTrank compounds recorded as approved after the 2023 label release, which
is more likely an artefact of name matching than a real approval. Applying the seven
corrections moves the ECE ratio by at most 0.31 and the reliability share by at most
2.2 percentage points, so the conclusions do not turn on them. Two of the three
post-2023 compounds are among those seven; the third, selumetinib, is not, because
its earliest matching record is also after the split boundary. The underlying
imprecision in the approval-year proxy remains.

Structure-based deduplication collapses salts and stereoisomers to a common parent,
which merged fourteen groups carrying conflicting labels. Removing them in full is
a defensible rule, fixed before its effect was examined, and its effect is not
one-directional --- re-running the decomposition on the retained pre-filter sets
gives 85.6\% against 86.1\% in DILIrank and 84.5\% against 82.5\% in DICTrank, a
rise in one and a fall in the other --- but it is still an intervention on the
analysis set, and the removed groups are listed in full so a reader can judge it.

ATC assignment carries two unresolved ambiguities: where KEGG lists several codes
for one drug the first alphabetically is taken, and where two drugs normalise to
the same name key the first encountered wins. Both would need the KEGG REST
interface, which was not reachable during this work. The class-L results are
therefore hypothesis-generating regardless of their nominal $p$-values.

The molecular representations are fingerprint- and descriptor-based. Graph neural
networks were not evaluated, on the judgement that a held-out set of 156 compounds
cannot distinguish architectures reliably; that judgement is itself a limitation
rather than a result.

The Murphy decomposition and the conformal analysis were computed for one
model--representation pair. ECE itself was checked across all 32 combinations, so
the calibration claim is not architecture-specific, but the decomposition and the
coverage analysis have not been repeated across the grid.

Two of the three candidate explanations were evaluated on AUROC only. The maturity
controls establish that censoring explains at most a quarter of the discrimination
gap in DILIrank and essentially none of it in DICTrank, and the similarity analysis
finds no significant relation between structural distance and per-compound accuracy,
but neither recomputes calibration. The claim that prevalence shift is the dominant
calibration term therefore rests on the nested prevalence controls alone. Repeating
the maturity and similarity analyses with ECE as the outcome would test it from a
second direction and is the most obvious next experiment.

Finally, the three-family structure used for multiple-comparison correction was
fixed after the analyses were run rather than pre-registered. The grouping criteria
are stated explicitly so that a reader can regroup the tests and recompute; the
confirmatory family survives at $q = 0.0070$ and the exploratory family yields
nothing significant under any grouping we examined.

\subsection{Implications for practice}

Three recommendations follow, in decreasing order of confidence.

Report calibration alongside discrimination whenever a temporal or prospective
split is used. In these data ECE was worse under the temporal split in all 32
model--representation combinations and at all three cut-off years, with the ratio
never falling below 1.6 across cut-offs and below 1.43 across the grid. Over the
same conditions $\Delta$AUROC varied by a factor of 2.4 across cut-offs in DILIrank,
swapped the rank of its 2005 and 2015 values between the two datasets, and changed
sign in one of the 32 combinations. The calibration signal was the more consistent of the two, and it
is the one usually omitted.

Do not treat scaffold splitting as evidence of prospective validity without
reporting the scaffold-size distribution. If most scaffolds are singletons, the
design is closer to a random split than its name suggests.

Treat an internally calibrated model as calibrated only for the period it was
calibrated on. Where a prospective use is intended, the relevant question is not
whether the probabilities were correct at the end of the training window but
whether anything is known about how the base rate will move afterwards.

\section{Conclusions}

Across two FDA regulatory toxicity benchmarks, splitting on approval year degrades
model performance in a way that AUROC substantially understates. The loss in
discrimination is 0.088 and 0.077 of AUROC; against a random split of the same size
expected calibration error rises from 0.131 to 0.438 and from 0.130 to 0.270,
ratios of 3.35 and 2.08, and the Murphy decomposition attributes 86.1\% and 82.5\%
of the increase in Brier score to the calibration term. The pattern held in all 32 combinations of four learners and four molecular
representations, across three cut-off years, ten random seeds and eight binning
schemes, and was confirmed independently by a distribution-free conformal coverage
analysis that requires no calibration metric at all.

Prevalence shift accounts for 62\% of the calibration collapse in DILIrank and 66\%
in DICTrank, and not for the rest. Chemical-space novelty and label censoring were tested against the discrimination
gap only. Chemical-space novelty accounts for none of it; label censoring accounts
for at most a quarter of it in DILIrank and essentially none in DICTrank. Whether
either contributes to the calibration collapse was not measured.
Bemis--Murcko scaffold splitting does not reproduce it, because in drug-like
reference sets most scaffolds occur once and the design reduces to something close
to a random split.

The practical consequence is that a model reported as losing "only" 0.09 of AUROC
under time-split validation may be producing probabilities that are wrong by a
factor that no reported number reveals. Calibration metrics are cheap to compute
and, in this setting, carry the more stable and more actionable signal. They should
be reported as a matter of course.
